\section{Positive replacement and exact rectangular refinement}\label{sec:matrixproof}
We prove Theorem~\ref{thm:matrix}. Every step uses the actual matrix profile
on the forward stack in \eqref{eq:actual}. In particular, positivity is
used to discard wrap terms only after replacing the signed profile.

\subsection{The spectral split and scalar inverse moment}
Set
\[
 \sigma=\alpha_D^{1/P}c,
 \qquad S_x=1_{(-\infty,\sigma]}(A(x))\R^r,
 \qquad N=\sum_x\dim S_x,\qquad I=\{x:S_x\ne\{0\}\}.
\]
Here $1_{(-\infty,\sigma]}(A(x))$ is the spectral projector. A degenerate
eigenspace at the cutoff is selected in full, so no choice of basis within
it enters the construction. Every nonpositive channel is selected. A positive
selected channel contributes exactly one to the capped budget, and an
unselected channel contributes its uncapped inverse power. Consequently
\begin{equation}\label{eq:cap-exact}
 F_c(A)=N+\alpha_Dc^P
                \sum_{x,j:\lambda_j(x)>\sigma}\lambda_j(x)^{-P},
 \qquad 0\le N\le F_c(A),\quad |I|\le N.
\end{equation}
Equality $\lambda=\sigma$ is included in $N$.

Put $T=\sigma+c(1+N)>\sigma$. Replace every selected eigenvalue by $T$,
leaving all unselected eigenvalues and every local spectral subspace
unchanged, and denote the resulting positive definite matrix by $A_T(x)$.
For $N=0$ this changes nothing. With
$a_T(x)=\lambda_{\min}(A_T(x))>0$, pointwise matrix positivity gives
\begin{equation}\label{eq:scalar-domination}
 Q_{A_T}(u)\ge Q_{a_T}(u).
\end{equation}
This controls the complete stacked vector, including all directional and
internal mixing. It requires no commutation between distinct matrices.
Moreover, the largest eigenvalue of $A_T(x)^{-P}$ is at most its trace, hence
\begin{align}\label{eq:replacement-moment}
 c^P\sum_xa_T(x)^{-P}
 &\le c^P\sum_x\tr(A_T(x)^{-P})\nonumber\\
 &=c^P\sum_{\lambda_j(x)>\sigma}\lambda_j(x)^{-P}
                  +N(c/T)^P\\
 &\le \frac{F_c(A)-N}{\alpha_D}+1.\nonumber
\end{align}
The last inequality uses $T\ge c(1+N)$ and
$N/(1+N)^P\le1$ for $P\ge1$. When $N=0$ the extra one is harmless.

\subsection{The scalar partition is derived, not assumed}
Theorem~\ref{thm:balanced223} in Appendix~\ref{sec:balanced223} applies to
every finite positive $a_T$, without an arithmetic-mean, lower-envelope or
upper-envelope condition. Its partition $\mathcal R_0$, combined with
\eqref{eq:scalar-domination}, satisfies
\begin{equation}\label{eq:positive-observation}
 c\normn{u}^2\le V^{-1}Q_{A_T}(u)+c\|C_{\mathcal R_0}u\|^2,
\end{equation}
by summing the same scalar inequality over all $h$ real components.
If $R_0=\#\mathcal R_0$, its exact count and
\eqref{eq:replacement-moment} yield
\begin{align}\label{eq:initial-count}
 R_0&\le2D\left[1+(2b+3)b(cG_D)^P\sum_xa_T(x)^{-P}\right]\nonumber\\
 &\le2D\left[1+B_D\left(1+\frac{F_c(A)-N}{\alpha_D}\right)\right].
\end{align}
There is no auxiliary local-gap or extension hypothesis in this application:
the scalar construction and its gap are proved in the appendix.

\subsection{Isolating every affected endpoint}
Let
\begin{equation}\label{eq:endpoints}
 Z=I\cup\bigcup_{i=1}^D(I+e_i),\qquad |Z|\le(D+1)|I|\le(D+1)N,
\end{equation}
where the translations are periodic. We refine $\mathcal R_0$ so that each
$z\in Z$ is a singleton member. Order $Z$ lexicographically in the fixed
chart. If $z$ is not already a singleton, let $R$ be its current rectangle.
Partition $R\setminus\{z\}$ according to the first coordinate in which a
point differs from $z$, and according to whether it is smaller or larger in
that coordinate. Earlier coordinates are fixed to those of $z$, later
coordinates range over their original intervals, and the differing coordinate
ranges below or above $z$. Each nonempty piece is an integer rectangle.
There are at most $2D$ such pieces. Adding $\{z\}$ replaces one rectangle
by at most $2D+1$, increasing the count by at most $2D$.
Earlier isolated points lie in disjoint singleton members and remain intact.
Thus the final partition $\mathcal R$ satisfies
\begin{equation}\label{eq:refined-count}
 R:=\#\mathcal R\le R_0+2D|Z|\le R_0+2D(D+1)N.
\end{equation}
A wrapping forward endpoint is simply its chart representative in $Z$;
this causes no alteration of the rectangle argument.

For each original parent $B$ and its final pieces $B_k$,
\[
 q_B(u)=\sum_k\sqrt{|B_k|/|B|}\,q_{B_k}(u),\qquad
 \sum_k|B_k|/|B|=1.
\]
Hilbert-space Cauchy with these exact weights proves
\begin{equation}\label{eq:refine-monotone}
 \|C_{\mathcal R_0}u\|^2\le\|C_{\mathcal R}u\|^2,
 \qquad V^{-1}\sum_{z\in Z}\|u(z)\|^2
                      \le\|C_{\mathcal R}u\|^2.
\end{equation}
The second assertion holds because each $z\in Z$ contributes the singleton
row $V^{-1/2}u(z)$. The remaining rows are still only component means on
rectangles; no gradient or nonrectangular observation has been introduced.

\subsection{Charging the complete signed correction}
Define the finite number
\begin{equation}\label{eq:negative-depth}
 L=\begin{cases}
 0,&N=0,\\
 \displaystyle\max_{x,j:\lambda_j(x)\le\sigma}(T-\lambda_j(x)),&N>0.
 \end{cases}
\end{equation}
The spectral construction gives
$0\le A_T(x)-A(x)\le L\Id$, with zero difference for $x\notin I$.
This is an operator inequality on all $r$ coordinates; in particular it
bounds all rotated off-diagonal contributions together. It does not assign
signs to individual cross terms. Therefore
\begin{align}\label{eq:correction}
 0\le Q_{A_T}(u)-Q_A(u)
 &\le L\sum_{x\in I}\|du(x)\|^2\nonumber\\
 &\le2L\sum_{x\in I}
      \left[D\|u(x)\|^2+\sum_i\|u(x+e_i)\|^2\right]\\
 &\le4DL\sum_{z\in Z}\|u(z)\|^2.\nonumber
\end{align}
Indeed each site has its own bad-start multiplicity at most $D$ and at most
$D$ incoming direction-labelled starts. All their endpoints belong to $Z$.
This counts parallel edges at period two separately, as required by the
actual form. At $n=1$ every difference vanishes and the estimate remains
valid. No periodic term has been dropped from the signed correction.

Insert \eqref{eq:correction} into \eqref{eq:positive-observation}, divide
the point sum by $V$, and use \eqref{eq:refine-monotone}. This proves
\eqref{eq:matrix-coercive} with $K=c+4DL$. The count is
\begin{align*}
 hR&\le2Dh\left[1+B_D+\frac{B_D}{\alpha_D}(F_c(A)-N)
                                      +(D+1)N\right]\\
 &\le2Dh\left[1+B_D+
                 \max\{B_D/\alpha_D,D+1\}F_c(A)\right],
\end{align*}
using the nonnegativity of both $N$ and $F_c(A)-N$. This completes the proof
of Theorem~\ref{thm:matrix}.

\begin{remark}[What the cap does and does not control]\label{rem:cap}
The only use of the budget after the local spectral split is the exact
identity \eqref{eq:cap-exact}. In particular its control of observation
number is compatible with arbitrarily negative selected eigenvalues, since
the penalty absorbs the actual finite correction. When $N=0$, the correction
vanishes and $K=c$. For a positive scalar profile the sharper direct scalar
construction also has $K=c$ without selecting or replacing any channels.
For completeness, if $A(x)\ge-M\Id$ for all $x$, then
$Q_A(u)\ge-4DM\sum_x\|u(x)\|^2$, because
$\sum_x\|du(x)\|^2\le4D\sum_x\|u(x)\|^2$.
For the singleton partition $\|C_{\mathcal R}u\|^2=\normn{u}^2$,
so $K=c+4DM$ proves its asserted finite-penalty alternative.
The matrix construction establishes an existence bound with explicit
constants, rather than a sharp spectral rank formula.
\end{remark}
